% !TeX root = ../main.tex
\section{End-to-End Audio-to-Tablature Models}
\label{sec:e2e}

End-to-end models predict string and fret directly from audio, so the
fingering can use timbre and the tablature constraints can shape the
transcription. Before deep learning, Barbancho et~al.~\cite{barbancho2012chords}
transcribed chords and their fingering from audio with an HMM over
fingering configurations.

\emph{Frame-wise classification.} TabCNN~\cite{wiggins2019tabcnn} applies a
small CNN to a constant-Q spectrogram and outputs, per frame, a distribution
over fret classes (open, frets, silent) for each of the six strings, trained
and evaluated on GuitarSet. It has no note-level decoding and no temporal
model. Cwitkowitz et~al.~\cite{cwitkowitz2022inhibition} add an inhibition loss
that discourages co-activation of note pairs that are unlikely on the
fretboard, estimated from symbolic tablature. FretNet~\cite{cwitkowitz2023fretnet}
estimates continuous-valued pitch contours grouped by string and fret, which
represents bends and vibrato that a semitone grid cannot.

\emph{Pre-training on synthetic data.} SynthTab~\cite{zang2024synthtab} shows
that pre-training a GTT model on DadaGP rendered through guitar plugins
improves results after fine-tuning and, more importantly, reduces overfitting
in cross-dataset tests; this is the main evidence that data scale, not
architecture, limits end-to-end GTT.

\emph{Electric guitar and techniques.} GOAT~\cite{loth2025goat} reports
preliminary results for a tablature transcription approach trained on its
DI-plus-amplifier data, and is the first large source of electric-guitar audio
paired with technique-annotated tablature.

\emph{Generative models.} Noise2Fret~\cite{simionato2026noise2fret} generates
tablature with a diffusion model over a continuous latent representation of
string and fret targets, conditioned on spectral features. Five auxiliary
losses (pitch-class distance, positional distance, circle-of-fifths distance,
string similarity, hand-span feasibility) encode playability; on GuitarSet and
GOAT the model outperforms its baselines and the losses give consistent gains.

\emph{Hybrid systems.} TART~\cite{gupta2025tart,gupta2026tart} is presented
as end-to-end but is a four-stage pipeline (audio-to-MIDI, technique
classifier, audio-conditioned fingering, tablature generator). Evaluated
zero-shot on GuitarSet, EGDB and noisy versions of both, it reports 81.35\%
audio-to-MIDI F50, 71.8\% string-fret Tab F1 and 54.08\% end-to-end Tab F1
on average. The gap between the last two numbers shows how errors compound
across stages.

All of these models take an isolated guitar as input. None is trained or
evaluated on guitar inside a band mix, and most are trained on GuitarSet,
which is acoustic.
